\section{系统总体设计}

\subsection{总体设计思路}

针对自然语言输入不稳定而 CAD 脚本要求稳定的矛盾，本文采用“前松后紧”的总体设计策略。即在输入端允许用户以自然语言自由描述需求，在中间层将需求收缩为严格受限的 Plan，在输出端只允许确定性模板生成脚本。其本质是通过中间表示把语言理解与几何生成解耦。

系统总体流程如下：读取用户需求文本；调用分析模块生成候选 JSON；解析并校验为 \texttt{AssemblyPlan}；根据零件形状调用对应渲染器生成局部脚本；将各零件脚本合并为完整 FeatureScript；最终通过 CLI 或 Web UI 向用户返回结果，并按需落盘保存。

\begin{figure}[t]
  \centering
  \resizebox{0.98\linewidth}{!}{\begin{tikzpicture}[
  node distance=0.8cm and 0.9cm,
  >=Latex,
  box/.style={draw, rounded corners, align=center, minimum width=3.0cm, minimum height=1.0cm, fill=blue!6},
  io/.style={draw, rounded corners, align=center, minimum width=2.6cm, minimum height=0.9cm, fill=green!8},
  store/.style={draw, rounded corners, align=center, minimum width=2.6cm, minimum height=0.9cm, fill=orange!10},
  line/.style={->, thick}
]

\node[io] (cli) {CLI\\Analyze / Generate / Build};
\node[io, right=of cli] (web) {Web UI\\Analyze / Generate / Full Build};
\node[box, below=of $(cli)!0.5!(web)$] (app) {Application Layer\\Use Cases / Service Orchestration};
\node[box, below left=of app] (analysis) {Requirement Analyzer\\Prompting / Provider / Parsing / Fallback};
\node[box, below=of app] (schema) {AssemblyPlan Schema\\Typed Validation};
\node[box, below right=of app] (gen) {FeatureScript Generator\\Deterministic Renderers / Merge};
\node[store, below=of analysis] (req) {Natural-Language\\Requirements};
\node[store, below=of schema] (plan) {Validated Plan\\JSON Artifact};
\node[store, below=of gen] (fs) {FeatureScript\\Artifact};

\draw[line] (cli) -- (app);
\draw[line] (web) -- (app);
\draw[line] (app) -- (analysis);
\draw[line] (app) -- (schema);
\draw[line] (app) -- (gen);
\draw[line] (req) -- (analysis);
\draw[line] (analysis) -- (schema);
\draw[line] (schema) -- (plan);
\draw[line] (plan) -- (gen);
\draw[line] (gen) -- (fs);

\end{tikzpicture}
}
  \caption{系统总体架构。系统围绕“需求分析、Plan 校验、确定性脚本生成”三层核心能力展开，并同时服务于 CLI 与 Web UI 两类入口。}
  \label{fig:architecture}
\end{figure}

\subsection{系统架构设计}

从代码结构来看，系统主要由 \texttt{analysis}、\texttt{models}、\texttt{generation}、\texttt{application}、\texttt{io}、\texttt{webui} 和 \texttt{cli} 七类模块组成，分别负责模型分析、Schema 定义、模板生成、流程编排、文件读写、页面服务和命令入口。该分层方式使 CLI 与 Web UI 不直接实现业务逻辑，而是通过应用层与服务层复用同一套分析和生成流程。

\subsection{数据结构设计}

系统的核心中间表示为 \texttt{AssemblyPlan}。该结构至少包含以下字段：
\begin{enumerate}[leftmargin=2em]
  \item \texttt{assembly\_name}：装配体标识；
  \item \texttt{description}：装配描述；
  \item \texttt{global\_params}：全局尺寸信息和坐标原点描述；
  \item \texttt{parts}：零件列表；
  \item \texttt{assembly\_relations}：零件之间的装配关系。
\end{enumerate}

其中每个零件需要描述 \texttt{id}、\texttt{name}、\texttt{shape}、\texttt{material\_hint}、\texttt{params}、\texttt{position} 和 \texttt{description}。当前系统仅支持五种基础形状，这一限制是为了确保模板映射明确、参数集合稳定。

\subsection{关键流程设计}

\subsubsection{Analyze 流程}

Analyze 流程用于将需求文本转换为 Plan。应用层首先读取需求文本，然后调用 \texttt{RequirementAnalyzer}。分析器会使用预设提示词向大模型发起请求，获取流式返回内容，再尝试解析为 JSON 对象，并通过模型层进行严格校验。若模型调用失败或返回异常，而当前输入又匹配仓库内置样例，则系统会启用示例兜底逻辑，保证演示链路可运行。

\subsubsection{Generate 流程}

Generate 流程面向已有 Plan 文件。系统读取并校验 Plan 后，遍历全部零件，根据零件形状选择对应渲染器生成 FeatureScript 代码片段。所有片段最终由合并模块封装为标准 FeatureScript 文件。

\subsubsection{Build 流程}

Build 是全流程入口。它既可以接收自然语言需求，也可以在给定 \texttt{--plan} 时跳过分析阶段，直接进入生成。最终系统会输出 Plan 文件和 FeatureScript 文件，适合用于教学演示和批量脚本化使用。

\subsection{交互设计}

\subsubsection{CLI 设计}

CLI 是系统的主入口，支持多子命令。其优势在于可脚本化、便于集成、适合开发者调试和自动化验证。命令参数允许用户指定输入文本、输入文件、输出目录、分析模型、API Key 等运行时配置。

\subsubsection{Web UI 设计}

Web UI 面向演示场景，采用本地 HTTP 服务加静态前端页面方式实现。页面支持需求输入、样例填充、Plan JSON 编辑、FeatureScript 预览、结果摘要展示以及日志回显。由于 Web UI 直接复用后端 \texttt{WebUIService} 的能力，因此其分析与生成结果与 CLI 保持一致。
